\documentclass[11pt,a4paper]{article}

\usepackage[utf8]{inputenc}
\usepackage[T1]{fontenc}
\usepackage{amsmath}
\usepackage{geometry}
\usepackage{hyperref}
\usepackage{enumitem}
\usepackage{booktabs}

\geometry{margin=2.5cm}
\hypersetup{
  colorlinks=true,
  linkcolor=blue,
  urlcolor=blue,
  citecolor=blue
}

\title{Supplementary Material: Reproducibility Details}
\author{St\'ephane Ga\"el R. Ekodeck, Serge Alain Ebele,\\
Prosper Rosaire Mama Assandje, and Ren\'e Ndoundam}
\date{}

\begin{document}
\maketitle

\noindent
This document supplements the article ``Boundary Information Geometry for
S-box Side-Channel Leakage: Walsh Covariance, Active Fisher Structure, and CPA
Belief Dynamics.''  It records the computational protocol used for the
numerical claims in the manuscript.  The computations are intentionally small
and exhaustive: for $s=4$, all input values, all key hypotheses, all
Walsh coordinates, and all active coordinate planes are enumerated.

\section{Software Environment}

The distributed scripts use Python 3 and standard numerical packages.  The core
dependencies are \texttt{numpy} for exact finite-array enumeration and linear
algebra, \texttt{scipy} for symmetric eigenspace computations, and
\texttt{matplotlib} for figure generation.  No randomized optimization,
training procedure, or stochastic sampling is used for the reported tables.
When an additional 4-bit bijection is used in \texttt{verify\_curvature.py},
the script uses a fixed table explicitly listed in the source code.

\section{Commands}

The main numerical outputs can be regenerated from the project root with:
\begin{verbatim}
python gen_figures.py
python compute_curvatures.py
python verify_curvature.py
\end{verbatim}

\noindent
The first command regenerates the Walsh-spectrum and local key-separation
figures.  The second command constructs the boundary covariance matrix,
projects onto the active eigenspace, computes the projected third cumulant
contractions, and evaluates the curvature diagnostic.  The third command is a
sanity check for the value $K=-1/4$ on PRESENT and on an additional fixed
random bijection.

\section{Arithmetic and Precision}

The Walsh coefficients, boundary covariance entries, and third cumulant entries
are obtained from finite sums over $\{0,1\}^4$.  These quantities are therefore
integer or rational before diagonalization.  The eigenspace projection and the
curvature table are evaluated in floating-point arithmetic after the exact
finite sums have been constructed.  The reported curvature values are
reproducible numerical diagnostics for the tested S-boxes under the projection
and normalization conventions used in the manuscript.

\section{Sign Conventions}

The Walsh transform convention is
\[
  \widehat S(\alpha,\beta)=\sum_{a\in\{0,1\}^s}
  (-1)^{\alpha\cdot a\oplus\beta\cdot S(a)}.
\]
The active-eigenspace curvature diagnostic follows the convention used in
Section 6 of the manuscript.  Changing the Riemann tensor sign convention would
change the sign of all reported sectional curvature values; the scripts and
the manuscript use the same convention.

\section{Reproducibility Coverage}

The computational protocol covers the quantities reported in the paper:
\begin{enumerate}[label=\normalfont(\roman*)]
  \item Walsh spectra and absolute-coefficient histograms;
  \item boundary covariance matrices and active ranks/eigenvalues;
  \item local Walsh/Fisher key-separation proxies;
  \item Hamming-weight leakage descriptors and pairwise leakage gaps;
  \item projected active-eigenspace curvature diagnostics for PRESENT, GIFT,
  and SKINNY.
\end{enumerate}

\section{Files}

\begin{center}
\begin{tabular}{ll}
\toprule
File & Role \\
\midrule
\texttt{gen\_figures.py} & Regenerates the manuscript figures. \\
\texttt{compute\_curvatures.py} & Computes active ranks, eigenvalues, and curvature values. \\
\texttt{verify\_curvature.py} & Performs an independent curvature sanity check. \\
\texttt{requirements.txt} & Lists the Python package dependencies. \\
\bottomrule
\end{tabular}
\end{center}

\end{document}
